%%--------------------------------------------------------------------
%1--------------------------------------------------------------------
\begin{glyph}{Command (令)}{command}\label{glyph:command}
Command.  Order.\\\hline
This 漢字 is also used as an honorific prefix in certain compounds.
\end{glyph}
\gindex{Command}{令}\strindex{5}{令}

\begin{comps}
\comprtwocone&亼 + 卩\\\hline
\comprthreecone&Assembly/Gathering + Seal/Stamp\\\hline
\end{comps}

\begin{remglyph}
\hooglig{Command} the \remcom{assembly} to give their \remcom{stamp} of approval.\\\hline
\end{remglyph}

\begin{prons}
\pronsrtwocone&\textbf{レイ (REI)}&---\\\hline
\pronsrthreecone&---&---\\\hline
\end{prons}

\begin{remprons}
\hooglig{Command} the \comon{rain} to stop.\\\hline
\end{remprons}

%{\middel}
%&\mbox{()}\\\hline
\begin{somme}
指令&finger + command = \textit{order; directive}&シ{\middel}レイ \mbox{(SHI{\middel}REI)}\\\hline
令和&command + harmony = \textit{Reiwa Era (2019/05/01---{\ldots})}&レイ{\middel}ワ \mbox{(REI{\middel}WA)}\\\hline
令達&command + achieve (reach) = \textit{delivered order}&レイタツ \mbox{(REI{\middel}TATSU)}\\\hline
令名&command + name (noted) = \textit{fame}&レイ{\middel}メイ \mbox{(REI{\middel}MEI)}\\\hline
令夫人&command + wife = \textit{Mrs; Lady; your wife}&レイ{\middel}フ{\middel}ジン \mbox{(REI{\middel}FU{\middel}JIN)}\\\hline
令兄&command + elder brother = \textit{your elder brother}&レイ{\middel}ケイ \mbox{(REI{\middel}KEI)}\\\hline
指令&finger + command = \textit{orders}&シ{\middel}レイ \mbox{(SHI{\middel}REI)}\\\hline
\notyet{号令}&title + command = \textit{command; order}&ゴウ{\middel}レイ \mbox{(G\={O}{\middel}REI)}\\\hline
\notyet{法令}&law + command = \textit{laws and ordinances}&ホウ{\middel}レイ \mbox{(H\={O}{\middel}REI)}\\\hline
\end{somme}

\begin{decomp}{6}
私&の&命令&は&絶対&だ。\\\hline
わたし&の&メイレイ&は&ゼッタイ&だ\\\hline
I&の&orders&は&absolute&be/is\\\hline
\multicolumn{6}{|r|}{\textit{My orders are absolute.}}\\\hline
\end{decomp}

\end{multicols}
%
%%--------------------------------------------------------------------
% command all climb
%2--------------------------------------------------------------------
\begin{multicols}{2}
\begin{glyph}{Measurement (尺)}{measurement}\label{glyph:measurement}
Measurement.  30.3 cm.
\end{glyph}
\gindex{Measurement}{尺}\strindex{4}{尺}

\begin{comps}
\comprtwocone&丶 + 尸\\\hline
\comprthreecone&Dot + Corpse/Body\\\hline
\end{comps}

\begin{remglyph}
\hooglig{Measure} the \remcom{dot} on the \remcom{corpse}.\\\hline
\end{remglyph}

\begin{prons}
\pronsrtwocone&\textbf{シャク (SHAKU)}&---\\\hline
\pronsrthreecone&---&---\\\hline
\end{prons}

\begin{remprons}
\hooglig{Measure} with a \comon{shakuhachi}.\\\hline
\end{remprons}

%{\middel}
%&\mbox{()}\\\hline
\begin{somme}
一尺&one + measurement = \textit{approx. 30 cm}&イッ{\middel}シャク \mbox{(IS{\middel}SHAKU)}\\\hline
尺八&measurement + eight = \textit{shakuhachi flute}&シャク{\middel}ハチ \mbox{(SHAKU{\middel}HACHI)}\\\hline
尺骨&30.3 cm + bone = \textit{Ulna}&シャッ{\middel}コツ \mbox{(SHAK{\middel}KOTSU)}\\\hline
尺地&measurement + ground = \textit{small land plot}&シャク{\middel}チ \mbox{(SHAKU{\middel}CHI)}\\\hline
尺寸&measurement + tiny = \textit{sth. tiny; trifle}&シャク{\middel}スン \mbox{(SHAKU{\middel}SUN)}\\\hline
\notyet{尺度}&measurement + degree = \textit{linear measure; scale}&シャク{\middel}ド \mbox{(SHAKU{\middel}DO)}\\\hline
\end{somme}

\begin{decomp}{7}
人&は&万物&の&尺度&で&ある。\\\hline
ひと&は&バンブツ&の&シャクド&で&ある\\\hline
person&は&creation&の&scale&で&ある\\\hline
\multicolumn{7}{|r|}{\textit{Man is the measure of all things.}}\\\hline
\end{decomp}

\end{multicols}
%
\pagebreak
%%--------------------------------------------------------------------
%3--------------------------------------------------------------------

\begin{multicols}{2}
\begin{glyph}{Flat (平)}{flat}\label{glyph:flat}
Flat.\\\hline
Extended meanings such as \textit{ordinary} and \textit{peaceful}.
\end{glyph}
\gindex{Flat}{平}\strindex{5}{平}

\begin{comps}
\comprtwocone&艹 + 干\\\hline
\comprthreecone&Grass/Herb/Plant + Dry\\\hline
\end{comps}

\begin{remglyph}
\hooglig{Flat}　\remcom{herbs} are \remcom{dry}.\\\hline
\end{remglyph}

\begin{prons}
\pronsrtwocone&\textbf{ヘイ (HEI)}&\textbf{ひら (hira)}\\\hline
\rye{3}{ヘイ is the more common of these readings.}
\pronsrthreecone&ビョウ (BY\={O})&たい (tai)\\\hline
\end{prons}

\begin{remprons}
\hooglig{The flats} \comkun{he rustled} went \comon{Haywire}.\\\hline
\hooglig{Flats} in \ucomkun{Tailand} will \ucomon{be yours}.\\\hline
\end{remprons}

%{\middel}
%&\mbox{()}\\\hline
\begin{somme}
平たい&\textit{flat; level}&ひら{\middel}たい \mbox{(hira{\middel}tai)}\\\hline
平和&flat + harmony = \textit{peace}&ヘイ{\middel}ワ \mbox{(HEI{\middel}WA)}\\\hline
太平洋&fat + flat + ocean = \textit{the Pacific Ocean}&タイ{\middel}ヘイ{\middel}ヨウ \mbox{(TAI{\middel}HEI{\middel}Y\={O})}\\\hline
平日&flat + sun = \textit{weekday}&ヘイ{\middel}ジツ \mbox{(HEI{\middel}JITSU)}\\\hline
平気&flat + spirit = \textit{all right; OK}&ヘイ{\middel}キ \mbox{(HEI{\middel}KI)}\\\hline
平服&flat + clothes = \textit{ordinary clothes}&ヘイ{\middel}フク \mbox{(HEI{\middel}FUKU)}\\\hline
平手&flat + hand = \textit{palm; open hand}&ひら{\middel}て \mbox{(hira{\middel}te)}\\\hline
公平&public + flat = \textit{fairness; impartiality}&コウ{\middel}ヘイ \mbox{(K\={O}{\middel}HEI)}\\\hline
\notyet{平和運動}&peace + exercise = \textit{peace movement}&ヘイ{\middel}ワ{\middel}ウン{\middel}ドウ \mbox{(HEI{\middel}WA{\middel}UN{\middel}D\={O})}\\\hline
\notyet{平泳ぎ}&flat + swim = \textit{breaststroke}&ひら{\middel}およ{\middel}ぎ \mbox{(hira{\middel}oyo{\middel}gi)}\\\hline
\notyet{地平線}&level ground + line = \textit{horizon}&チ{\middel}ヘイ{\middel}セン \mbox{(CHI{\middel}HEI{\middel}SEN)}\\\hline
\end{somme}

\begin{decomp}{5}
平和&は&とても&大切&です。\\\hline
ヘイワ&は&とても&タイセツ&です\\\hline
peace&は&very&important&be/is\\\hline
\multicolumn{5}{|r|}{\textit{Peace is of great importance.}}\\\hline
\end{decomp}

\end{multicols}
%%--------------------------------------------------------------------
%4--------------------------------------------------------------------
\begin{multicols}{2}
\begin{glyph}{Degree (度)}{degree}\label{glyph:degree}
Degree.  Extent.\\\hline
This 漢字 can also mean \textit{times} in the sense of \textit{how many times}.
\end{glyph}
\gindex{Degree}{度}\strindex{9}{度}

\begin{comps}
\comprtwocone&广 +　廿 + 又\\\hline
\comprthreecone&Slanted　Roof + Twenty　+ Again\\\hline
\end{comps}

\begin{remglyph}
\hooglig{The extent} of the \remcom{slanted roof} making a \remcom{20} \hooglig{degree} angle \remcom{again}.\\\hline
\end{remglyph}

\begin{prons}
\pronsrtwocone&\textbf{ド (DO)}&\textbf{たび (tabi)}\\\hline
\pronsrthreecone&タク、ト (TAKU, TO)&---\\\hline
\end{prons}

\begin{remprons}
\hooglig{The extent} of \comkun{tubby} \comon{dots} \ucomon{tucked} at the \ucomon{top}.\\\hline
\end{remprons}

%{\middel}
%&\mbox{()}\\\hline
\begin{somme}
\splitter{度する}{tr}&\textit{to redeem (from sin)}&ド{\middel}する \mbox{(DO{\middel}suru)}\\\hline
度合い&degree + match = \textit{degree; extent}&ド{\middel}あ{\middel}い \mbox{(DO{\middel}a{\middel}i)}\\\hline
度々&degree + degree = \textit{often; over and over}&たび{\middel}たび \mbox{(tabi{\middel}tabi)}\\\hline
一度&one + degree (times) = \textit{once; one time}&イチ{\middel}ド \mbox{(ICHI{\middel}DO)}\\\hline
丁度&block + degree = \textit{exactly}&チョウ{\middel}ド \mbox{(CH\={O}{\middel}DO)}\\\hline
{\diff}度外れ&degree + outside = \textit{extraordinary; excessive}&ド{\wsplit}はず{\middel}れ \mbox{(DO{\wsplit}hazu{\middel}re)}\\\hline
\notyet{度数}&degree + number = \textit{frequency; incidence}&ド{\middel}スウ \mbox{(DO{\middel}S\={U})}\\\hline
\notyet{度し難い}&degree + difficult = \textit{beyond help}&ど{\middel}し{\wsplit}がた{\middel}い \mbox{(do{\middel}shi{\wsplit}gata{\middel}i)}\\\hline
\notyet{制度}&regulation + degree = \textit{system}&セイ{\middel}ド \mbox{(SEI{\middel}DO)}\\\hline
\notyet{温度}&warm + degree = \textit{temperature}&オン{\middel}ド \mbox{(ON{\middel}DO)}\\\hline
\end{somme}

\begin{decomp}{2}
さあ、&もう一度。\\\hline
さあ&もうイチド\\\hline
come now&once more\\\hline
\multicolumn{2}{|r|}{\textit{Come on, try again.}}\\\hline
\end{decomp}
\end{multicols}
\pagebreak
%%--------------------------------------------------------------------
%5--------------------------------------------------------------------

\begin{multicols}{2}
\begin{glyph}{Edge (際)}{edge}\label{glyph:edge}
Edge.  Limit.\\\hline
A secondary meaning has to do with \textit{occasion}.
\end{glyph}
\gindex{Edge}{際}\strindex{14}{際}

\begin{comps}
\comprtwocone&阝阜 + 祭\\\hline
\comprthreecone&Hill/Mound + Festival (p. \pageref{glyph:festival})\\\hline
\end{comps}

\begin{remglyph}
Edging the hill we found a festival.\\\hline
\end{remglyph}

\begin{prons}
\pronsrtwocone&\textbf{サイ (SAI)}&\textbf{きわ (kiwa)}\\\hline
\pronsrthreecone&---&---\\\hline
\end{prons}

\begin{remprons}
\hooglig{Edgy} \comon{psycho} \comkun{key-warden}.\\\hline
\end{remprons}

%{\middel}
%&\mbox{()}\\\hline
\begin{somme}
際どい&\textit{risky; questionable}&きわ{\middel}どい \mbox{(kiwa{\middel}doi)}\\\hline
際して&\textit{when}&サイ{\middel}して \mbox{(SAI{\middel}shite)}\\\hline
際会&edge + meet = \textit{meeting; facing}&サイ{\middel}カイ \mbox{(SAI{\middel}KAI)}\\\hline
水際&water + edge = \textit{beach; waterside}&みず{\middel}ぎわ \mbox{(mizu{\middel}giwa)}\\\hline
国際&country + edge = \textit{international}&コク{\middel}サイ \mbox{(KOKU{\middel}SAI)}\\\hline
際立つ&edge + stand = \textit{to be prominent}&きわ{\wsplit}だ{\middel}つ \mbox{(kiwa{\middel}da{\wsplit}tsu)}\\\hline
実際&actual + edge = \textit{in practice; in actuality}&ジッ{\middel}サイ \mbox{(JIS{\middel}SAI)}\\\hline
際物&edge + thing = \textit{seasonal goods}&きわ{\middel}もの \mbox{(kiwa{\middel}mono)}\\\hline
交際&mix + edge = \textit{company}&コウ{\middel}サイ \mbox{(K\={O}{\middel}SAI)}\\\hline
手際&hand + edge = \textit{performance}&て{\middel}ぎわ \mbox{(te{\middel}giwa)}\\\hline
一際&one + edge = \textit{conspicuously}&ひと{\middel}きわ \mbox{(hito{\middel}kiwa)}\\\hline
\end{somme}

\begin{decomp}{5}
英語&は&国際語&に&成った。\\\hline
エイゴ&は&コクサイゴ&に&なった\\\hline
English&は&international language&に&to become\\\hline
\multicolumn{5}{|r|}{\textit{English has become an international language.}}\\\hline
\end{decomp}

\end{multicols}

%%--------------------------------------------------------------------
%6--------------------------------------------------------------------
\begin{multicols}{2}
\begin{glyph}{Lend (貸)}{lend}\label{glyph:lend}
Lend.
\end{glyph}
\gindex{Lend}{貸}\strindex{12}{貸}

\begin{comps}
\comprtwocone&代 + 貝\\\hline
\comprthreecone&Substitute (p. \pageref{glyph:substitute}) + Shellfish\\\hline
\end{comps}

\begin{remglyph}
\hooglig{Lending} \remcom{shellfish} to the \remcom{substitute}.\\\hline
\end{remglyph}

\begin{prons}
\pronsrtwocone&---&\textbf{か (ka)}\\\hline
\rye{3}{The 訓読み is frequently voiced outside the first position.}
\pronsrthreecone&タイ (TAI)&---\\\hline
\end{prons}

\begin{remprons}
\hooglig{Lending} to \comkun{customers} is really a \ucomon{tight} spot.\\\hline
\end{remprons}

%{\middel}
%&\mbox{()}\\\hline
\begin{somme}
\splitter{貸す}{tr}&\textit{to lend/loan}&か{\middel}す \mbox{(ka{\middel}su)}\\\hline
\splitter{貸し出す}{tr}&lend + exit = \textit{to lend/loan}&か{\middel}し{\wsplit}だ{\middel}す \mbox{(ka{\middel}shi{\wsplit}da{\middel}su)}\\\hline
貸付&lend + attach = \textit{loan}&だし{\middel}つけ \mbox{(kashi{\middel}tsuke)}\\\hline
貸切&lend + cut = \textit{reserved; chartered}&かし{\middel}きり \mbox{(kashi{\middel}kiri)}\\\hline
貸出&lend + exit = \textit{lending; loaning}&かし{\middel}だし \mbox{(kashi{\middel}dashi)}\\\hline
貸間&lend + interval = \textit{room to let}&かし{\middel}ま \mbox{(kashi{\middel}ma)}\\\hline
貸主&lend + primary = \textit{lender; landlord}&かし{\middel}ぬし \mbox{(kashi{\middel}nushi)}\\\hline
貸家&lend + house = \textit{house for rent}&かし{\middel}や \mbox{(kashi{\middel}ya)}\\\hline
貸元&lend + basis = \textit{financier; lender}&かし{\middel}もと \mbox{(kashi{\middel}moto)}\\\hline
貸本屋&lend + bookshop = \textit{rental library}&かし{\middel}ホン{\middel}や \mbox{(kashi{\middel}HON{\middel}ya)}\\\hline
\end{somme}

\begin{decomp}{6}
金&を&貸して&友&を&失え。\\\hline
キン&を&かして&とも&を&うしなえ\\\hline
gold&を&to lend&friend&を&to lose\\\hline
\multicolumn{6}{|r|}{\textit{Lend your money and lose your friend.}}\\\hline
\end{decomp}

\end{multicols}
\pagebreak
%%--------------------------------------------------------------------
%7--------------------------------------------------------------------
\begin{multicols}{2}
\begin{glyph}{Taste (味)}{taste}\label{glyph:taste}
Taste.  Flavour.
\end{glyph}
\gindex{Taste}{味}\strindex{8}{味}

\begin{comps}
\comprtwocone&口 + 未\\\hline
\comprthreecone&Mouth + Not Yet (p. \pageref{glyph:not_yet})\\\hline
\end{comps}

\begin{remglyph}
\hooglig{The tastes} of \remcom{vampires} have \remcom{not yet} been approved.\\\hline
\end{remglyph}

\begin{prons}
\pronsrtwocone&\textbf{ミ (MI)}&\textbf{あじ (aji)}\\\hline
\pronsrthreecone&---&---\\\hline
\end{prons}

\begin{remprons}
\hooglig{Tasting} \comon{mitigated} by \comkun{ajin}.\\\hline
{\warn}\textit{Ajin} are demi-humans in the `Ajin' アニメ series.\\\hline
\end{remprons}

%{\middel}
%&\mbox{()}\\\hline
\begin{somme}
\splitter{味わう}{tr}&\textit{to taste (sth)}&あじ{\middel}わう \mbox{(aji{\middel}wau)}\\\hline
\splitter{味がする}{tr}&\textit{to taste of (sth)}&あじ{\middel}がする \mbox{(aji{\middel}gasuru)}\\\hline
\splitter{味わい知る}{tr}&\textit{to taste and know (sth)}&あじ{\middel}わい{\middel}し{\middel}る \mbox{(aji{\middel}wai{\middel}shi{\middel}ru)}\\\hline
味な&\textit{smart; clever; witty}&あじ{\middel}な \mbox{(aji{\middel}na)}\\\hline
気味&spirit + taste = \textit{feeling/sensation}&キ{\middel}ミ \mbox{(KI{\middel}MI)}\\\hline
味気ない&taste + spirit = \textit{wearisome; dull; wretched; vain}&あじ{\middel}ケ{\middel}ない \mbox{(aji{\middel}KE{\middel}nai)}\\\hline
意味&mind + taste = \textit{meaning}&イ{\middel}ミ \mbox{(I{\middel}MI)}\\\hline
味方&taste + direction = \textit{friend; ally}&ミ{\middel}かた \mbox{(MI{\middel}kata)}\\\hline
\end{somme}

\begin{irrr}
美味しい&beautiful + taste = \textit{delicious}&おい{\middel}しい \mbox{(oi{\middel}shii)}\\\hline
不味い&not + taste = \textit{bad taste; awful}&ま{\middel}ず{\middel}い \mbox{(ma{\middel}zu{\middel}i)}\\\hline
\end{irrr}

\begin{decomp}{5}
喫煙&は&自殺&を&意味する。\\\hline
キツエン&は&ジサツ&を&イミする\\\hline
smoking&は&suicide&を&meaning\\\hline
\multicolumn{5}{|r|}{\textit{Smoking means suicide.}}\\\hline
\end{decomp}

\end{multicols}
%\begin{decomp}{9}
%彼女&の&その&問題&の&説明&は&無意味&だった。\\\hline
%かのジョ&の&その&モンダイ&の&セツメイ&は&ムイミ&だった\\\hline
%she/her&の&that&problem&の&explanation&は&meaningless&be/is\\\hline
%\multicolumn{9}{|r|}{\textit{Her explanation of the problem added up to the nonsense.}}\\\hline
%\end{decomp}

%%--------------------------------------------------------------------
%8--------------------------------------------------------------------
\begin{multicols}{2}
\begin{glyph}{Field (野)}{field}\label{glyph:field}
Field.  The wild.\\\hline
Interestingly, this 漢字 is used to denote a \textit{political party in opposition}.
\end{glyph}
\gindex{Field}{野}\strindex{11}{野}

\begin{comps}
\comprtwocone&里 + 予\\\hline
\comprthreecone&Hamlet + Beforehand (p. \pageref{glyph:beforehand})\\\hline
\end{comps}

\begin{remglyph}
\hooglig{Fields} in the \remcom{hamlet} must be watered \remcom{beforehand}.\\\hline
\end{remglyph}

\begin{prons}
\pronsrtwocone&\textbf{ヤ (YA)}&\textbf{の (no)}\\\hline
\pronsrthreecone&---&---\\\hline
\end{prons}

\begin{remprons}
\hooglig{In the wild} \comon{yachting} with the \comkun{nobs}.\\\hline
{\warn}\textit{Nob} is a person of wealth or high social position.\\\hline
\end{remprons}

%{\middel}
%&\mbox{()}\\\hline
\begin{somme}
野心&field + heart = \textit{ambition}&ヤ{\middel}シン \mbox{(YA{\middel}SHIN)}\\\hline
分野&part + field = \textit{field (of study, etc.)}&ブン{\middel}ヤ \mbox{(BUN{\middel}YA)}\\\hline
野山&field + mountain = \textit{hills and fields}&の{\middel}やま \mbox{(no{\middel}yama)}\\\hline
野生動物&field + life + animal = \textit{wild animal}&ヤ{\middel}セイ{\middel}ドウ{\middel}ブツ \mbox{(YA{\middel}SEI{\middel}D\={O}{\middel}BUTSU)}\\\hline
平野&flat + field = \textit{open field}&ヘイ{\middel}ヤ \mbox{(HEI{\middel}YA)}\\\hline
野生&field + nature = \textit{living in the wild}&ヤ{\middel}セイ \mbox{(YA{\middel}SEI)}\\\hline
大野&large + field = \textit{large field}&おお{\middel}の \mbox{(\={o}{\middel}no)}\\\hline
野外&field + outside = \textit{outdoors}&ヤ{\middel}ガイ \mbox{(YA{\middel}GAI)}\\\hline
\notyet{野原}&field + original = \textit{field(s)}&no{\middel}hara \mbox{(の{\middel}はら)}\\\hline
\notyet{野菜}&field + vegetable = \textit{vegetable}&ヤ{\middel}サイ \mbox{(YA{\middel}SAI)}\\\hline
\end{somme}

\begin{decomp}{6}
彼&は&大変&な&野心家&だった。\\\hline
かれ&は&タイヘン&な&ヤシンカ&だった\\\hline
he&は&immense&な&ambitious person&だった\\\hline
\multicolumn{6}{|r|}{\textit{He was a man of great ambition.}}\\\hline
\end{decomp}

\end{multicols}
\pagebreak
%%--------------------------------------------------------------------
%9--------------------------------------------------------------------

\begin{multicols}{2}
\begin{glyph}{Wide (広)}{wide}\label{glyph:wide}
Wide.  Broad.  Spacious.
\end{glyph}
\gindex{Wide}{広}\strindex{5}{広}

\begin{comps}
\comprtwocone&广 + 厶\\\hline
\comprthreecone&Slanted roof + Self/Private\\\hline
\end{comps}

\begin{remglyph}
\hooglig{Wide} \remcom{slanted roofs} are for \remcom{private} use, \hooglig{broadly} speaking.\\\hline
\end{remglyph}

\begin{prons}
\pronsrtwocone&\textbf{コウ (K\={O})}&\textbf{ひろ (hiro)}\\\hline
\pronsrthreecone&---&---\\\hline
\end{prons}

\begin{remprons}
\hooglig{Wide} open, \comkun{Hiroshima} was its own \comon{cause} of defeat.\\\hline
\end{remprons}

%{\middel}
%&\mbox{()}\\\hline
\begin{somme}
\splitter{広がる}{intr}&\textit{to spread (out)}&ひろ{\middel}がる \mbox{(hiro{\middel}garu)}\\\hline
\splitter{広げる}{tr}&\textit{to spread (sth)}&ひろ{\middel}げる \mbox{(hiro{\middel}geru)}\\\hline
広がり&\textit{expanse; extent}&ひろ{\middel}がり \mbox{(hiro{\middel}gari)}\\\hline
広さ&\textit{extent}&ひろ{\middel}さ \mbox{(hiro{\middel}sa)}\\\hline
広間&wide + interval = \textit{hall; saloon}&ひろ{\middel}ま \mbox{(hiro{\middel}ma)}\\\hline
広い&\textit{wide; spacious}&ひろ{\middel}い \mbox{(hiro{\middel}i)}\\\hline
広島&wide + island = \textit{Hiroshima}&ひろ{\middel}しま \mbox{(hiro{\middel}shima)}\\\hline
広大&wide + large = \textit{extensive; immense}&コウ{\middel}ダイ \mbox{(K\={O}{\middel}DAI)}\\\hline
広々&wide $(\times 2)$ = \textit{extensive; spacious}&ひろ{\middel}びろ \mbox{(hiro{\middel}biro)}\\\hline
広言&wide + speak = \textit{boasting}&コウ{\middel}ゲン \mbox{(K\={O}{\middel}GEN)}\\\hline
\notyet{広場}&wide + place = \textit{public square}&ひろ{\middel}ば \mbox{(hiro{\middel}ba)}\\\hline
\end{somme}

\begin{decomp}{4}
その&川&は&広い。\\\hline
その&かわ&は&ひろい\\\hline
that/the&river&は&wide\\\hline
\multicolumn{4}{|r|}{\textit{The river is wide.}}\\\hline
\end{decomp}

\end{multicols}

%%--------------------------------------------------------------------
%10-------------------------------------------------------------------
\begin{multicols}{2}
\begin{glyph}{Stone (石)}{stone}\label{glyph:stone}
Stone.
\end{glyph}
\gindex{Stone}{石}\strindex{5}{石}

\begin{comps}
\comprtwocone&石\\\hline
\comprthreecone&Stone/Rock\\\hline
\end{comps}

\begin{remglyph}
Part of the 漢字 radicals (\#{112}).\\\hline
\end{remglyph}

\begin{prons}
\pronsrtwocone&\textbf{セキ (SEKI)}&\textbf{いし (ishi)}\\\hline
\rye{3}{セキ is more common outside of the first position.}
\pronsrthreecone&シャク、コク (SHAKU, KOKU)&---\\\hline
\end{prons}

\begin{remprons}
\hooglig{The stone}-colored \comkun{Ishihara} test that \ucomon{cocked} my interest was \ucomon{shockingly} \comon{sexy}.\\\hline
\end{remprons}

%{\middel}
%&\mbox{()}\\\hline
\begin{somme}
石&\textit{stone}&いし \mbox{(ishi)}\\\hline
小石&small + stone = \textit{pebble}&こ{\middel}いし \mbox{(ko{\middel}ishi)}\\\hline
化石&change + stone = \textit{fossil}&カ{\middel}セキ \mbox{(KA{\middel}SEKI)}\\\hline
一石二鳥&one + stone + two + bird = \textit{kill two birds with one stone}&イッ{\middel}セキ{\middel}ニ{\middel}チョウ \mbox{(IS{\middel}SEKI{\middel}NI{\middel}CH\={O})}\\\hline
石英&stone + brilliance = \textit{quartz}&セキ{\middel}エイ \mbox{(SEKI{\middel}EI)}\\\hline
石頭&stone + head = \textit{obstinate person; stubbornness}&いし{\middel}あたま \mbox{(ishi{\middel}atama)}\\\hline
大理石&large + reason + stone = \textit{marble}&ダイ{\middel}リ{\middel}セキ \mbox{(DAI{\middel}RI{\middel}SEKI)}\\\hline
定石&fixed + stone = \textit{standard practice}&ジョウ{\middel}セキ \mbox{(J\={O}{\middel}SEKI)}\\\hline
石川県&stone + river + prefecture = \textit{Ishikawa prefecture}&いし{\middel}かわ{\middel}ケン \mbox{(ishi{\middel}kawa{\middel}KEN)}\\\hline
\notyet{石室}&stone + room = \textit{stone hut; tomb}&セキ{\middel}シツ \mbox{(SEKI{\middel}SHITSU)}\\\hline
\notyet{試金石}&assay + stone = \textit{touchstone}&シ{\middel}キン{\middel}セキ \mbox{(SHI{\middel}KIN{\middel}SEKI)}\\\hline
\end{somme}

\begin{decomp}{4}
橋&は&石造り&だ。\\\hline
はし&は&いしづくり&だ\\\hline
bridge&は&made of stone&だ\\\hline
\multicolumn{4}{|r|}{\textit{The bridge is made of stone.}}\\\hline
\end{decomp}

\end{multicols}

\pagebreak
%%--------------------------------------------------------------------
%11-------------------------------------------------------------------

\begin{multicols}{2}
\begin{glyph}{China (漢)}{china}\label{glyph:china}
China.  The Han Dynasty.
\end{glyph}
\gindex{China}{漢}\strindex{13}{漢}

\begin{comps}
\comprtwocone&水氵氺 + 艹 + 口 + 夫\\\hline
\comprthreecone&Water + Grass + Mouth + Husband (p. \pageref{glyph:husband})\\\hline
\end{comps}

\begin{remglyph}
\hooglig{Chinese} \remcom{water} \remcom{herbs} are for \remcom{vampire} \remcom{husbands}.\\\hline
\end{remglyph}

\begin{prons}
\pronsrtwocone&\textbf{カン (KAN)}&---\\\hline
\pronsrthreecone&---&---\\\hline
\end{prons}

\begin{remprons}
\hooglig{The Chinese} are \comon{cunning}.\\\hline
\end{remprons}

%{\middel}
%&\mbox{()}\\\hline
\begin{somme}
漢語&China + word = \textit{Sino-Japanese word}&カン{\middel}ゴ \mbox{(KAN{\middel}GO)}\\\hline
漢方&China + direction = \textit{traditional Chinese medicine}&カン{\middel}ポウ \mbox{(KAN{\middel}P\={O})}\\\hline
漢方薬&China + direction + medicine = \textit{Chinese herbal medicine}&カン{\middel}ポウ{\middel}ヤク \mbox{(KAN{\middel}P\={O}{\middel}YAKU)}\\\hline
漢和&China + harmony = \textit{China and Japan}&カン{\middel}ワ \mbox{(KAN{\middel}WA)}\\\hline
漢民族&China + people + tribe = \textit{Chinese people}&カン{\middel}ミン{\middel}ゾク \mbox{(KAN{\middel}MIN{\middel}ZOKU)}\\\hline
漢人&China + person = \textit{Chinese person}&カン{\middel}ジン \mbox{(KAN{\middel}JIN)}\\\hline
漢土&China + earth = \textit{China}&カン{\middel}ド \mbox{(KAN{\middel}DO)}\\\hline
漢文&China + literature = \textit{Chinese classical writing}&カン{\middel}ブン \mbox{(KAN{\middel}BUN)}\\\hline
\notyet{漢字試験}&Kanji + test = \textit{Kanji test}&カン{\middel}ジ{\middel}シ{\middel}ケン \mbox{(KAN{\middel}JI{\middel}SHI{\middel}KEN)}\\\hline
\end{somme}
\end{multicols}

\begin{decomp}{7}
この&漢字&は&どういう&意味&です&か。\\\hline
この&カンジ&は&どういう&イミ&です&か\\\hline
this&Kanji&は&what kind of&meaning&be/is&か\\\hline
\multicolumn{7}{|r|}{\textit{What does this kanji mean?}}\\\hline
\end{decomp}

%%--------------------------------------------------------------------
%12-------------------------------------------------------------------
\begin{multicols}{2}
\begin{glyph}{Different (別)}{different}\label{glyph:different}
Different.\\\hline
The ideas of \textit{separation} and \textit{farewell}.
\end{glyph}
\gindex{Different}{別}\strindex{7}{別}

\begin{comps}
\comprtwocone&口 + 勹 + 刂刁刀ク\\\hline
\comprthreecone&Mouth + Embracing/Wrap + Knife/Sword\\\hline
\end{comps}

\begin{remglyph}
\hooglig{Different} \remcom{vampires} \remcom{embraced} the idea of a \remcom{sword} \hooglig{farewell}.\\\hline
\end{remglyph}

\begin{prons}
\pronsrtwocone&\textbf{ベツ (BETSU)}&\textbf{わか (waka)}\\\hline
\pronsrthreecone&---&わ (wa)\\\hline
\end{prons}

\begin{remprons}
\hooglig{Different} \ucomkun{wanderers} made \comkun{wakaberry} \comon{bets}.\\\hline
\end{remprons}

%{\middel}
%&\mbox{()}\\\hline
\begin{somme}
別れ&\textit{parting; separation}&わか{\middel}れ \mbox{(waka{\middel}re)}\\\hline
別に&\textit{(not) particularly}&ベツ{\middel}に \mbox{(BETSU{\middel}ni)}\\\hline
\splitter{別れる}{intr}&\textit{to part \ldots}&わか{\middel}れる \mbox{(waka{\middel}reru)}\\\hline
別々&different $(\times 2)$ = \textit{separate; individual}&ベツ{\middel}ベツ \mbox{(BETSU{\middel}BETSU)}\\\hline
送別&send + different = \textit{farewell; send-off}&ソウ{\middel}ベツ \mbox{(S\={O}{\middel}BETSU)}\\\hline
別名&different + name = \textit{alias; pen name}&ベツ{\middel}メイ \mbox{(BETSU{\middel}MEI)}\\\hline
別料金&different + fee = \textit{extra cost/charge}&ベツ{\middel}リョウ{\middel}キン \mbox{(BETSU{\middel}RY\={O}{\middel}KIN)}\\\hline
\notyet{特別}&special + different = \textit{special; extraordinary}&トク{\middel}ベツ \mbox{(TOKU{\middel}BETSU)}\\\hline
\notyet{区別}&partition + different = \textit{distinction}&ク{\middel}ベツ \mbox{(KU{\middel}BETSU)}\\\hline
\notyet{選別}&select + different = \textit{selection}&セン{\middel}ベツ \mbox{(SEN{\middel}BETSU)}\\\hline
\notyet{別館}&different + manor = \textit{annex}&ベッ{\middel}カン \mbox{(BEK{\middel}KAN)}\\\hline
\end{somme}

\begin{decomp}{3}
別々&に&払います。\\\hline
ベツベツ&に&はらいます\\\hline
separately&に&to pay\\\hline
\multicolumn{3}{|r|}{\textit{Please bill us separately.}}\\\hline
\end{decomp}
\end{multicols}

\pagebreak
%%--------------------------------------------------------------------
%13-------------------------------------------------------------------
\begin{multicols}{2}
\begin{glyph}{Garden (園)}{garden}\label{glyph:garden}
Garden.
\end{glyph}
\gindex{Garden}{園}\strindex{13}{園}

\begin{comps}
\comprtwocone&囗 + (袁 = 土 + 口 + 衣衤)\\\hline
\comprthreecone&Prison + (Long Kimono = Earth + Mouth + Clothes)\\\hline
\end{comps}

\begin{remglyph}
\hooglig{The garden} was a \remcom{prison} of \remcom{long kimonos}.\\\hline
\end{remglyph}

\begin{prons}
\pronsrtwocone&\textbf{エン (EN)}&---\\\hline
\pronsrthreecone&---&その (sono)\\\hline
\end{prons}

\begin{remprons}
\hooglig{Garden} \comon{entrance} graced with \ucomkun{sonorous} sounds.\\\hline
{\warn}\textit{Sonorous} means imposingly deep and full (of a sound).\\\hline
\end{remprons}

%{\middel}
%&\mbox{()}\\\hline
\begin{somme}
園&\textit{garden (esp. man-made)}&その \mbox{(sono)}\\\hline
園生&garden + life = \textit{garden (esp. with trees)}&エン{\middel}セイ \mbox{(EN{\middel}SEI)}\\\hline
公園&public + garden = \textit{park}&コウ{\middel}エン \mbox{(K\={O}{\middel}EN)}\\\hline
学園&study + garden = \textit{educational institute}&ガク{\middel}エン \mbox{(GAKU{\middel}EN)}\\\hline
動物園&animal + garden = \textit{zoo}&ドウ{\middel}ブツ{\middel}エン \mbox{(D\={O}{\middel}BUTSU{\middel}EN)}\\\hline
園内&garden + inside = \textit{inside the garden}&エン{\middel}ナイ \mbox{(EN{\middel}NAI)}\\\hline
園遊会&garden + play + meet = \textit{garden party}&エン{\middel}ユウ{\middel}カイ \mbox{EN{\middel}Y\={U}{\middel}KAI}\\\hline
園丁&garden + block = \textit{gardener}&エン{\middel}テイ \mbox{(EN{\middel}TEI)}\\\hline
\notyet{園長}&garen + long = \textit{head of a garden}&エン{\middel}チョウ \mbox{(EN{\middel}CH\={O})}\\\hline
\notyet{園庭}&garden + yard = \textit{garden}&エン{\middel}テイ \mbox{(EN{\middel}TEI)}\\\hline
\notyet{園池}&garden + pond = garden with a pond&エン{\middel}チ \mbox{(EN{\middel}CHI)}\\\hline
\end{somme}

\begin{decomp}{3}
公園&で&降りました。\\\hline
コウエン&で&おりました\\\hline
park&で&descend\\\hline
\multicolumn{3}{|r|}{\textit{I got off at the park.}}\\\hline
\end{decomp}

\end{multicols}

%%--------------------------------------------------------------------
%14-------------------------------------------------------------------
\begin{multicols}{2}
\begin{glyph}{All (皆)}{all}\label{glyph:all}
All.
\end{glyph}
\gindex{All}{皆}\strindex{9}{皆}

\begin{comps}
\comprtwocone&比 + 白\\\hline
\comprthreecone&Compare (p. \pageref{glyph:compare}) + White\\\hline
\end{comps}

\begin{remglyph}
\hooglig{All} \remcom{comparisons} of different tones of \remcom{white}.\\\hline
\end{remglyph}

\begin{prons}
\pronsrtwocone&---&\textbf{みな (mina)}\\\hline
\rye{3}{This reading is often pronounced \mbox{みんな} for emphasis.}
\pronsrthreecone&カイ (KAI)&---\\\hline
\end{prons}

\begin{remprons}
\hooglig{All} \comkun{minatory} \ucomon{kites}.\\\hline
{\warn}\textit{Minatory} expresses or conveys a threat.\\\hline
\end{remprons}

%{\middel}
%&\mbox{()}\\\hline
\begin{somme}
皆&\textit{all; everybody}&みな \mbox{(mina)}\\\hline
皆々&all $(\times 2)$ = \textit{everyone}&みな{\middel}みな \mbox{(mina{\middel}mina)}\\\hline
皆して&\textit{everyone at once}&みんな{\middel}して \mbox{(minna{\middel}shite)}\\\hline
皆目&all + eye = \textit{entirely; (not) at all}&カイ{\middel}モク \mbox{(KAI{\middel}MOKU)}\\\hline
皆無&all + nothingness = \textit{nil; bugger-all}&カイ{\middel}ム \mbox{(KAI{\middel}MU)}\\\hline
皆伝&all + transmit = \textit{initiation into a discipline}&カイ{\middel}デン \mbox{(KAI{\middel}DEN)}\\\hline
会具&all + tool = \textit{set (of clothes, gear, \ldots)}&カイ{\middel}グ \mbox{(KAI{\middel}GU)}\\\hline
皆中&all + middle = \textit{hitting all the targets (archery)}&カイ{\middel}チュウ \mbox{(KAI{\middel}CH\={U})}\\\hline
一切皆空&everything + all + empty = \textit{all is vanity}&いっ{\middel}さい{\middel}かい{\middel}くう\mbox{(IS{\middel}SAI{\middel}KAI{\middel}K\={U})}\\\hline
\end{somme}

\begin{decomp}{3}
皆&幸せ&だ。\\\hline
みんな&しあわせ&だ\\\hline
everyone&happiness&be/is\\\hline
\multicolumn{3}{|r|}{\textit{All are happy.}}\\\hline
\end{decomp}

\end{multicols}

\pagebreak
%%--------------------------------------------------------------------
%15-------------------------------------------------------------------
\begin{multicols}{2}
\begin{glyph}{Immediate (即)}{immediate}\label{glyph:immediate}
Immediate.  At once.
\end{glyph}
\gindex{Immediate}{即}\strindex{7}{即}

\begin{comps}
\comprtwocone&艮 + 卩\\\hline
\comprthreecone&Boundary + Seal/Stamp\\\hline
\end{comps}

\begin{remglyph}
\hooglig{Immediate} use of \remcom{good} \remcom{stamps} on the envelope's \remcom{boundary}.\\\hline
\end{remglyph}

\begin{prons}
\pronsrtwocone&\textbf{ソク (SOKU)}&---\\\hline
\pronsrthreecone&---&---\\\hline
\end{prons}

\begin{remprons}
\hooglig{Immediately}, the \comon{sokkie} came to an end.\\\hline
\end{remprons}

%{\middel}
%&\mbox{()}\\\hline
\begin{somme}
即日&immediate + sun (day) = \textit{on the same day}&ソク{\middel}ジツ \mbox{(SOKU{\middel}JITSU)}\\\hline
即夜&immediate + night = \textit{on the same night}&ソク{\middel}ヤ \mbox{(SOKU{\middel}YA)}\\\hline
即死&immediate + death = \textit{instant death}&ソク{\middel}シ \mbox{(SOKU{\middel}SHI)}\\\hline
即売&immediate + sell = \textit{sale on the spot}&ソク{\middel}バイ \mbox{(SOKU{\middel}BAI)}\\\hline
即金&immediate + gold = \textit{cash down}&ソッ{\middel}キン \mbox{(SOK{\middel}KIN)}\\\hline
即物的&immediate + matter + target = \textit{matter-of-fact; practical}&ソク{\middel}ブツ{\middel}テキ (SOKU{\middel}BUTSU{\middel}TEKI)\\\hline
即急&immediate + hasten = \textit{quick; sudden}&ソッ{\middel}キュウ \mbox{(SOK{\middel}KY\={U})}\\\hline
即今&immediate + today = \textit{at the moment}&ソッ{\middel}コン \mbox{(SOK{\middel}KON)}\\\hline
\notyet{即日速達}&same day + express delivery = \textit{same-day special delivery}&ソク{\middel}ジツ{\middel}ソク{\middel}タツ (SOKU{\middel}JITSU{\middel}SOKU{\middel}TATSU)\\\hline
\end{somme}

\begin{decomp}{4}
その&犬&は&即死した。\\\hline
その&いぬ&は&ソクシした\\\hline
that&dog&は&instant death\\\hline
\multicolumn{4}{|r|}{\textit{The dog was killed on the spot.}}\\\hline
\end{decomp}
\end{multicols}
%%--------------------------------------------------------------------
%16-------------------------------------------------------------------
\begin{multicols}{2}
\begin{glyph}{Shop (店)}{shop}\label{glyph:shop}
Shop.  Store.
\end{glyph}
\gindex{Shop}{店}\strindex{8}{店}

\begin{comps}
\comprtwocone&广 + 占\\\hline
\comprthreecone&Slanted Roof + Occupy\\\hline
\end{comps}

\begin{remglyph}
\hooglig{Shops} under \remcom{slanted roofs} are \remcom{occupied}.\\\hline
\end{remglyph}

\begin{prons}
\pronsrtwocone&\textbf{テン (TEN)}&\textbf{みせ (mise)}\\\hline
\rye{3}{テン is the more common reading.}
\pronsrthreecone&---&---\\\hline
\end{prons}

\begin{remprons}
\hooglig{The shop} was filled with \comon{tense} \comkun{misery}.\\\hline
\end{remprons}

%{\middel}
%&\mbox{()}\\\hline
\begin{somme}
店主&shop + primary = \textit{shopkeeper}&テン{\middel}シュ \mbox{(TEN{\middel}SHU)}\\\hline
店屋&shop + roof = \textit{shop; store}&みせ{\middel}や \mbox{(mise{\middel}ya)}\\\hline
開店&open + shop = \textit{shop/store opening}&カイ{\middel}テン \mbox{(KAI{\middel}TEN)}\\\hline
店開き&shop + open = \textit{starting a business}&みせ{\wsplit}びら{\middel}き \mbox{(mise{\wsplit}bira{\middel}ki)}\\\hline
書店&write + shop = \textit{bookstore}&ショ{\middel}テン \mbox{(SHO{\middel}TEN)}\\\hline
店先&shop + precede = \textit{store front}&みせ{\middel}さき \mbox{(mise{\middel}saki)}\\\hline
店頭&shop + head = \textit{over-the-counter; store front}&テン{\middel}トウ \mbox{(TEN{\middel}T\={O})}\\\hline
売店&sell + shop = \textit{stand; kiosk}&バイ{\middel}テン \mbox{(BAI{\middel}TEN)}\\\hline
店の上がり&shop + upper = \textit{income from a shop}&みせ{\wsplit}の{\wsplit}あ{\middel}がり \mbox{(mise{\wsplit}no{\wsplit}a{\middel}gari)}\\\hline
店屋物&shop + thing = \textit{take-out (food)}&テン{\middel}や{\middel}もの　\mbox{(TEN{\middel}ya{\middel}mono)}\\\hline
\end{somme}

\begin{decomp}{5}
書店&で&手&に&入ります。\\\hline
ショテン&で&て&に&はいります\\\hline
bookstore&で&hand&に&to get in\\\hline
\multicolumn{5}{|r|}{\textit{You can get it at a bookstore.}}\\\hline
\end{decomp}

\end{multicols}
\pagebreak
%%--------------------------------------------------------------------
%17-------------------------------------------------------------------
\begin{multicols}{2}
\begin{glyph}{Institution (院)}{institution}\label{glyph:institution}
Institution, in the sense of \textit{public buildings} such as hospitals, or \textit{bodies} such as legislatures.
\end{glyph}
\gindex{Institution}{院}\strindex{10}{院}

\begin{comps}
\comprtwocone&阝阜 + 完\\\hline
\comprthreecone&Hill/Mound + Perfect\\\hline
\end{comps}

\begin{remglyph}
\hooglig{The institution} on the \remcom{hill} is \remcom{perfect}.\\\hline
\end{remglyph}

\begin{prons}
\pronsrtwocone&\textbf{イン (IN)}&---\\\hline
\pronsrthreecone&---&---\\\hline
\end{prons}

\begin{remprons}
\hooglig{Institution} \comon{ink}.\\\hline
\end{remprons}

%{\middel}
%&\mbox{()}\\\hline
\begin{somme}
寺院&temple + institution = \textit{Buddhist temple}&ジ{\middel}イン \mbox{(JI{\middel}IN)}\\\hline
入院&enter + institution = \textit{to be admitted to hospital}&ニュウ{\middel}イン \mbox{(NY\={U}{\middel}IN)}\\\hline
院本&institution + main = \textit{drama; playbook}&イン{\middel}ポン \mbox{(IN{\middel}PON)}\\\hline
学院&study + institution = \textit{academy}&ガク{\middel}イン \mbox{(GAKU{\middel}IN)}\\\hline
院生&institution + life = \textit{graduate student}&イン{\middel}セイ \mbox{(IN{\middel}SEI)}\\\hline
上院&upper + institution = \textit{upper house}&ジョウ{\middel}イン \mbox{(J\={O}{\middel}IN)}\\\hline
\notyet{院卒}&institution + graduate = \textit{grad-school graduate}&イン{\middel}ソツ \mbox{(IN{\middel}SOTSU)}\\\hline
\notyet{病院}&illness + institution = \textit{hospital}&ビョウ{\middel}イン \mbox{(BY\={O}{\middel}IN)}\\\hline
\notyet{院試}&institution + test = \textit{grad-school entrance exam}&イン{\middel}シ \mbox{(IN{\middel}SHI)}\\\hline
\end{somme}

\begin{decomp}{5}
病院&は&ここ&から&近い。\\\hline
ビョウイン&は&ここ&から&ちかい\\\hline
hospital&は&here&from&near\\\hline
\multicolumn{5}{|r|}{\textit{The hospital is near here.}}\\\hline
\end{decomp}

\end{multicols}
%%--------------------------------------------------------------------
%18-------------------------------------------------------------------
\begin{multicols}{2}
\begin{glyph}{Climb (登)}{climb}\label{glyph:climb}
Climb.\\\hline
An important secondary meaning has to do with \textit{attendance}.
\end{glyph}
\gindex{Climb}{登}\strindex{12}{登}

\begin{comps}
\comprtwocone&癶 + 豆\\\hline
\comprthreecone&Footsteps/Departure + Bean\\\hline
\end{comps}

\begin{remglyph}
\hooglig{Climb} from your \remcom{departure} into a \remcom{kind of vessel}.\\\hline
\end{remglyph}

\begin{prons}
\pronsrtwocone&\textbf{トウ (T\={O})}&---\\\hline
\pronsrthreecone&ト (TO)&のぼ (nobo)\\\hline
\rye{3}{The 訓読み for this 漢字 serves as an alternate form of \mbox{上る} (のぼ{\middel}る/nobo{\middel}ru), which we met on p. \pageref{glyph:upper}.}
\end{prons}

\begin{remprons}
\hooglig{Climbing} is \comon{torture} when there is \ucomkun{nobody} at the \ucomon{top} to help you.\\\hline
\end{remprons}

%{\middel}
%&\mbox{()}\\\hline
\begin{somme}
登校&climb + school = \textit{attend school}&トウ{\middel}コウ \mbox{(T\={O}{\middel}K\={O})}\\\hline
登り口&climb + mouth = \textit{base (of mountain)}&のぼ{\middel}り{\wsplit}ぐち \mbox{(nobo{\middel}ri{\wsplit}guchi)}\\\hline
%集団登校&group + attend school = \textit{going to school in groups}&シュウ{\middel}ダン{\middel}トウ{\middel}コウ \mbox{(SH\={U}{\middel}DAN{\middel}T\={O}{\middel}K\={O})}\\\hline
登山&climb + mountain = \textit{mountain climbing}&ト{\middel}ザン \mbox{(TO{\middel}ZAN)}\\\hline
登記&climb + note = \textit{registry; registration}&トウ{\middel}キ \mbox{(T\={O}{\middel}KI)}\\\hline
登校日&attend school + day = \textit{school day}&トウ{\middel}コウ{\middel}び \mbox{(T\={O}{\middel}K\={O}{\middel}bi)}\\\hline
登記所&registry + location = \textit{registry office}&トウ{\middel}キ{\middel}ショ \mbox{(T\={O}{\middel}KI{\middel}SHO)}\\\hline
木登り&tree + climb = \textit{tree climbing}&き{\wsplit}のぼ{\middel}り \mbox{(ki{\wsplit}nobo{\middel}ri)}\\\hline
\notyet{登場}&climb + place = \textit{make an entrance}&トウ{\middel}ジョウ \mbox{(T\={O}{\middel}J\={O})}\\\hline
\notyet{登場人物}&entrance + character = \textit{character (in a play/novel)}&トウ{\middel}ジョウ{\middel}ジン{\middel}ブツ \mbox{(T\={O}{\middel}J\={O}{\middel}JIN{\middel}BUTSU)}\\\hline
{\iscov}登竜門&climb + dragon + gate = \textit{gateway to success}&トウ{\middel}リュウ{\middel}モン \mbox{(T\={O}{\middel}RY\={U}{\middel}MON)}\\\hline
\end{somme}

\begin{decomp}{3}
登校する&所&です。\\\hline
トウコウする&ところ&です\\\hline
going to school&place&be/is\\\hline
\multicolumn{3}{|r|}{\textit{I am going to school.}}\\\hline
\end{decomp}

\end{multicols}
\pagebreak
%%--------------------------------------------------------------------
%19-------------------------------------------------------------------
\begin{multicols}{2}
\begin{glyph}{Original (原)}{original}\label{glyph:original}
Original: things in their \textit{original} or \textit{primitive state}.\\\hline
Many compounds with this character also convey the idea of a \textit{plain} or \textit{field} (i.e. an original place that is \textit{unspoiled} or \textit{wilderness-like}).
\end{glyph}
\gindex{Original}{原}\strindex{10}{原}

\begin{comps}
\comprtwocone&厂 + 白 + 小\\\hline
\comprthreecone&Cliff + White + Small\\\hline
\end{comps}

\begin{remglyph}
\hooglig{Original} \remcom{cliffs} are \remcom{white} and \remcom{small}.\\\hline
\end{remglyph}

\begin{prons}
\pronsrtwocone&\textbf{ゲン (GEN)}&\textbf{はら (hara)}\\\hline
\pronsrthreecone&---&---\\\hline
\end{prons}

\begin{remprons}
\hooglig{Originally}, \comon{Genghis-Khan} committed \comkun{hara-kiri}.\\\hline
\end{remprons}

%{\middel}
%&\mbox{()}\\\hline
\begin{somme}
野原&field + original = \textit{field(s)}&の{\middel}はら \mbox{(no{\middel}hara)}\\\hline
高原&tall + original = \textit{plateau; tableland}&コウ{\middel}ゲン \mbox{(K\={O}{\middel}GEN)}\\\hline
原子力&original + child + strength = \textit{atomic energy; nuclear power}&ゲン{\middel}シ{\middel}リョク \mbox{(GEN{\middel}SHI{\middel}RYOKU)}\\\hline
原理&original + reason = \textit{fundamental truth}&ゲン{\middel}リ \mbox{(GEN{\middel}RI)}\\\hline
原書&original + write = \textit{original document}&ゲン{\middel}ショ \mbox{(GEN{\middel}SHO)}\\\hline
原住民&original + reside + people = \textit{indigenous people}&ゲン{\middel}ジュウ{\middel}ミン \mbox{(GEN{\middel}J\={U}{\middel}MIN)}\\\hline
\notyet{原始}&original + start = \textit{origin; genesis}&ゲン{\middel}シ \mbox{(GEN{\middel}SHI)}\\\hline
\notyet{原料}&original + fee = \textit{raw materials}&ゲン{\middel}リョウ \mbox{(GEN{\middel}RY\={O})}\\\hline
\notyet{原作}&original + make = \textit{original work}&ゲン{\middel}サク \mbox{(GEN{\middel}SAKU)}\\\hline
\end{somme}

\begin{irrr}
川原&river + original = \textit{dry riverbed}&かわ{\middel}ら \mbox{(kawa{\middel}ra)}\\\hline
\end{irrr}
\end{multicols}

\begin{decomp}{8}
野原&に&は&六頭&の&羊&が&いた。\\\hline
のはら&に&は&ロクトウ&の&ひつじ&が&いた\\\hline
field&に&は&six heads&の&sheep&が&to be\\\hline
\multicolumn{8}{|r|}{\textit{There were six sheep in the field.}}\\\hline
\end{decomp}
%%--------------------------------------------------------------------
%20-------------------------------------------------------------------
\begin{multicols}{2}
\begin{glyph}{Late (晩)}{late}\label{glyph:late}
Late.  Evening.\\\hline
This 漢字 also imparts the idea of \textit{late} in a number of more poetic words, e.g. 晩年.
\end{glyph}
\gindex{Late}{晩}\strindex{12}{晩}

\begin{comps}
\comprtwocone&日 + 免\\\hline
\comprthreecone&Sun/Day + Exemption\\\hline
\end{comps}

\begin{remglyph}
\hooglig{Late at evening}, the \remcom{sun} is \remcom{exempted}.\\\hline
\end{remglyph}

\begin{prons}
\pronsrtwocone&\textbf{バン (BAN)}&---\\\hline
\pronsrthreecone&---&---\\\hline
\end{prons}

\begin{remprons}
\hooglig{Late} \comon{bunnies}.\\\hline
\end{remprons}

%{\middel}
%&\mbox{()}\\\hline
\begin{somme}
晩年&late + year = \textit{one's later years}&バン{\middel}ネン \mbox{(BAN{\middel}NEN)}\\\hline
晩婚化&late marriage + change = \textit{increase in avg. marriage age}&バン{\middel}コン{\middel}カ \mbox{(BAN{\middel}KON{\middel}KA)}\\\hline
晩方&late + direction = \textit{toward evening}&バン{\middel}がた \mbox{(BAN{\middel}gata)}\\\hline
晩食&late + eat = \textit{supper}&バン{\middel}ショク \mbox{(BAN{\middel}SHOKU)}\\\hline
毎晩&every + late = \textit{every evening}&マイ{\middel}バン \mbox{(MAI{\middel}BAN)}\\\hline
昨晩&past + late = \textit{last night}&サク{\middel}バン \mbox{(SAKU{\middel}BAN)}\\\hline
一晩中&one night + middle = \textit{all night long}&ひと{\middel}バン{\middel}ジュウ \mbox{(hito{\middel}BAN{\middel}J\={U})}\\\hline
\notyet{今晩は}&now + late = \textit{``good evening''}&コン{\middel}バン{\middel}は \mbox{(KON{\middel}BAN{\middel}wa)}\\\hline
\notyet{晩成}&late + become = \textit{late completion/bloomer}&バン{\middel}セイ \mbox{(BAN{\middel}SEI)}\\\hline
\end{somme}

\begin{decomp}{4}
今晩&暇&か&な？\\\hline
コンバン&ひま&か&な\\\hline
tonight&spare time&か&な\\\hline
\multicolumn{4}{|r|}{\textit{Are you free tonight?}}\\\hline
\end{decomp}

\end{multicols}
\pagebreak
%%--------------------------------------------------------------------
%21-------------------------------------------------------------------
\begin{multicols}{2}
\begin{glyph}{Entwine (絡)}{entwine}\label{glyph:entwine}
Entwine.  Become entangled.
\end{glyph}
\gindex{Entwine}{絡}\strindex{12}{絡}

\begin{comps}
\comprtwocone&糸 + 各\\\hline
\comprthreecone&Thread + Each (p. \pageref{glyph:each})\\\hline
\end{comps}

\begin{remglyph}
\hooglig{Entwine} (the) \remcom{threads} \remcom{each} time.\\\hline
\end{remglyph}

\begin{prons}
\pronsrtwocone&\textbf{ラク (RAKU)}&\textbf{から (kara)}\\\hline
\pronsrthreecone&---&---\\\hline
\end{prons}

\begin{remprons}
\hooglig{Entangled} with \comkun{caramel} \comon{luckily}.\\\hline
\end{remprons}

%{\middel}
%&\mbox{()}\\\hline
\begin{somme}
絡げ&\textit{bundle; sheaf}&から{\middel}げ \mbox{(kara{\middel}ge)}\\\hline
絡み&\textit{linkage; entanglement}&から{\middel}み \mbox{(kara{\middel}mi)}\\\hline
\splitter{絡む}{intr}&\textit{to entwine; become entangled}&から{\middel}む\mbox{(kara{\middel}mu)}\\\hline
\splitter{絡み合う}{intr}&entwine + match = \textit{to get tangled (together)}&から{\middel}み{\wsplit}あ{\middel}う \mbox{(kara{\middel}mi{\wsplit}a{\middel}u)}\\\hline
\splitter{絡み付く}{intr}&entwine + attach = \textit{cling to; coil around}&から{\middel}み{\wsplit}つ{\middel}く \mbox{(kara{\middel}mi{\wsplit}tsu{\middel}ku)}\\\hline
絡子&entwine + child = \textit{Zen monk's waistcoat}&ラク{\middel}ス \mbox{(RAKU{\middel}SU)}\\\hline
頭絡&head + entwine = \textit{bridle}&トウ{\middel}ラク \mbox{(T\={O}{\middel}RAKU)}\\\hline
短絡&short + entwine = \textit{short circuit}&タン{\middel}ラク \mbox{(TAN{\middel}RAKU)}\\\hline
絡み酒&entangle + Sak\'{e} = \textit{a drunken interaction}&から{\middel}み{\wsplit}ざけ \mbox{(kara{\middel}mi{\wsplit}zake)}\\\hline
絡め取る&entangle + take = \textit{to apprehend}&から{\middel}め{\wsplit}と{\middel}る \mbox{(kara{\middel}me{\wsplit}to{\middel}ru)}\\\hline
\end{somme}
\end{multicols}

\begin{decomp}{6}
此の&つる草&は&木&に&絡みつきます。\\\hline
この&つるソウ&は&き&に&からみつきます\\\hline
this&vine&は&tree&に&to twine around\\\hline
\multicolumn{6}{|r|}{\textit{This vine winds around trees.}}\\\hline
\end{decomp}
%%--------------------------------------------------------------------
%22-------------------------------------------------------------------
\begin{multicols}{2}
\begin{glyph}{Gather (集)}{gather}\label{glyph:gather}
Gather.  Assemble.
\end{glyph}
\gindex{Gather}{集}\strindex{12}{集}

\begin{comps}
\comprtwocone&隹 + 木\\\hline
\comprthreecone&Old birds + Tree\\\hline
\end{comps}

\begin{remglyph}
\hooglig{Gathering} \remcom{old birds} in the \remcom{tree}.\\\hline
\end{remglyph}

\begin{prons}
\pronsrtwocone&\textbf{シュウ (SH\={U})}&\textbf{あつ (atsu)}\\\hline
\pronsrthreecone&---&つど (tsudo)\\\hline
\end{prons}

\begin{remprons}
\hooglig{Gather} at the \comkun{atsumori} play with our \comon{shoes} on and \ucomkun{stew dough}.\\\hline
{\warn}The \textit{Atsumori} play focuses on a young samurai killed in the Genpei War.\\\hline
\end{remprons}

%{\middel}
%&\mbox{()}\\\hline
\begin{somme}
\splitter{集まる}{intr}&\textit{gather; come together}&あつ{\middel}まる \mbox{(atsu{\middel}maru)}\\\hline
\splitter{集める}{tr}&\textit{gather; collect (sth)}&あつ{\middel}める \mbox{(atsu{\middel}meru)}\\\hline
集まり&\textit{gathering; assembly}&あつ{\middel}まり \mbox{(atsu{\middel}mari)}\\\hline
全集&complete + gather = \textit{complete works/collection}&ゼン{\middel}シュウ \mbox{(ZEN{\middel}SH\={U})}\\\hline
集中&gather + middle = \textit{concentration}&シュウ{\middel}チュウ \mbox{(SH\={U}{\middel}CH\={U})}\\\hline
集計&gather + measure = \textit{aggregation}&シュウ{\middel}ケイ \mbox{(SH\={U}{\middel}KEI)}\\\hline
集結&gather + bind = \textit{massing (of troops)}&シュウ{\middel}ケツ \mbox{(SH\={U}{\middel}KETSU)}\\\hline
集団&gather + group = \textit{group; mass}&シュウ{\middel}ダン \mbox{(SH\={U}{\middel}DAN)}\\\hline
集合場所&gathering + place = \textit{meeting place}&シュウ{\middel}ゴウ{\middel}ば{\middel}ショ \mbox{(SH\={U}{\middel}G\={O}{\middel}ba{\middel}SHO)}\\\hline
集配人&collection \&delivery + person = \textit{postman}&シュウ{\middel}ハイ{\middel}ジン \mbox{(SH\={U}{\middel}HAI{\middel}JIN)}\\\hline
\end{somme}

\begin{decomp}{3}
集中力&が&ありません。\\\hline
しゅうちゅうりょく&が&ありません\\\hline
concentration&が&\sout{to come about}\\\hline
\multicolumn{3}{|r|}{\textit{I have difficulty concentrating}}\\\hline
\end{decomp}

\end{multicols}
\pagebreak
%%--------------------------------------------------------------------
%23-------------------------------------------------------------------
\begin{multicols}{2}
\begin{glyph}{Type (式)}{type}\label{glyph:type}
The primary meaning is that of a \textit{type} or \textit{style} of something.\\\hline
It is also very commonly used, however, in the sense of \textit{ceremony}.
\end{glyph}
\gindex{Type}{式}\strindex{6}{式}

\begin{comps}
\comprtwocone&弋 + 工\\\hline
\comprthreecone&Shoot/Ceremony + Work/Craft\\\hline
\end{comps}

\begin{remglyph}
\hooglig{Type} of \remcom{shooting ceremony} for prosperous \remcom{work and crafts}.\\\hline
\end{remglyph}

\begin{prons}
\pronsrtwocone&\textbf{シキ (SHIKI)}&---\\\hline
\pronsrthreecone&---&---\\\hline
\end{prons}

\begin{remprons}
\hooglig{Type} of \comon{Sheikh}.\\\hline
{\warn}\textit{Sheikh} refers to a religious Muslim scholar.\\\hline
\end{remprons}

%{\middel}
%&\mbox{()}\\\hline
\begin{somme}
形式&form + type = \textit{form; formality}&ケイ{\middel}シキ \mbox{(KEI{\middel}SHIKI)}\\\hline
自動式&self + move + type = \textit{automatic}&ジ{\middel}ドウ{\middel}シキ \mbox{(JI{\middel}D\={O}{\middel}SHIKI)}\\\hline
式次第&ceremony + order = \textit{program of a ceremony}&シキ{\middel}シ{\middel}ダイ \mbox{(SHIKI{\middel}SHI{\middel}DAI)}\\\hline
式辞&ceremony + phrase = \textit{ceremonial address}&シキ{\middel}ジ \mbox{(SHIKI{\middel}JI)}\\\hline
式神&ceremony + god = \textit{form of magic / divination}&シキ{\middel}がみ \mbox{(SHIKI{\middel}gami)}\\\hline
式法&ceremony + law = \textit{ceremony; manners}&シキ{\middel}ホウ \mbox{(SHIKI{\middel}H\={O})}\\\hline
式年祭&ceremony + anniversary = \textit{imperial memorial ceremony}&シキ{\middel}ネン{\middel}サイ \mbox{(SHIKI{\middel}NEN{\middel}SAI)}\\\hline
\notyet{成人式}&become + person + type = \textit{Coming of Age ceremony}&セイ{\middel}ジン{\middel}セキ \mbox{(SEI{\middel}JIN{\middel}SHIKI)}\\\hline
\end{somme}

\begin{decomp}{5}
朝食&は&バイキング&形式&だった。\\\hline
チョウショク&は&バイキング&ケイシキ&だった\\\hline
breakfast&は&smorgasbord&form&だった\\\hline
\multicolumn{5}{|r|}{\textit{Breakfast is a smorgasbord.}}\\\hline
\end{decomp}

\end{multicols}
%%--------------------------------------------------------------------
%24-------------------------------------------------------------------
\begin{multicols}{2}
\begin{glyph}{Passage (通)}{passage}\label{glyph:passage}
The general concept is of \textit{passing through and along}.\\\hline
This is a very versatile 漢字 that can mean a variety of things, from \textit{street} (it often appears on road signs), to the ideas of \textit{thoroughly getting to know a subject} or \textit{having things go according to plan}.
\end{glyph}
\gindex{Passage}{通}\strindex{10}{通}

\begin{comps}
\comprtwocone&マ + 用 + 辶\\\hline
\comprthreecone&MA + Utilize + Road/Walk\\\hline
\end{comps}

\begin{remglyph}
\hooglig{Passages} for \remcom{mothers-to-be} {use} a {road} to be \remcom{walked}.\\\hline
\end{remglyph}

\begin{prons}
\pronsrtwocone&\textbf{ツウ (TS\={U})}&\textbf{とお (t\={o})}\\\hline
\rye{3}{ツウ is by far the more common reading in the first position.\\\hline
とお is never voiced when forming verbs.}
\pronsrthreecone&ツ (TSU)&かよ (kayo)\\\hline
\end{prons}

\begin{remprons}
\hooglig{The passage} to \comon{Tsuu'tina} was \comkun{tormenting} because of the \ucomon{tsunami} \comkun{knockout}.\\\hline
{\warn}\textit{Tsuu'Tina} is an indigenous group in Alberta, Canada.\\\hline
\end{remprons}

%{\middel}
%&\mbox{()}\\\hline
\begin{somme}
通る (intr.)&\textit{go along; pass by}&とお{\middel}る \mbox{(t\={o}{\middel}ru)}\\\hline
通す (tr.)&\textit{pierce; let through}&とお{\middel}す \mbox{(t\={o}{\middel}su)}\\\hline
大通り&large + passage = \textit{a main street}&おお　どお{\middel}り \mbox{(\={o} d\={o}{\middel}ri)}\\\hline
交通&mix + passage = \textit{traffic; communication}&コウ{\middel}ツウ \mbox{(K\={O}{\middel}TS\={U})}\\\hline
通学&passage + study = \textit{school commute}&ツウ{\middel}ガク \mbox{(TS\={U}{\middel}GAKU)}\\\hline
通知&passage + know = \textit{notice; notification}&ツウ{\middel}チ \mbox{(TS\={U}{\middel}CHI)}\\\hline
通行止め&traffic passage + stop = \textit{road closure}&ツウ{\middel}コウ{\middel}ど{\middel}め \mbox{(TS\={U}{\middel}K\={O}{\middel}do{\middel}me)}\\\hline
\end{somme}

\begin{decomp}{3}
交通量&が&少なかった。\\\hline
コウツウリョウ&が&すくなかった\\\hline
traffic&が&few\\\hline
\multicolumn{3}{|r|}{\textit{There wasn't much traffic.}}\\\hline
\end{decomp}

\end{multicols}
\pagebreak
%%--------------------------------------------------------------------
%25-------------------------------------------------------------------

\begin{multicols}{2}
\begin{glyph}{Past (昨)}{past}\label{glyph:past}
Past.
\end{glyph}
\gindex{Past}{昨}\strindex{9}{昨}

\begin{comps}
\comprtwocone&日 + 乍\\\hline
\comprthreecone&Sun + Although/While\\\hline
\end{comps}

\begin{remglyph}
\hooglig{The past}, \remcom{though}, is clear as \remcom{day}.\\\hline
\end{remglyph}

\begin{prons}
\pronsrtwocone&\textbf{サク (SAKU)}&---\\\hline
\pronsrthreecone&---&---\\\hline
\end{prons}

\begin{remprons}
\hooglig{The past} really \comon{sucks}.\\\hline
\end{remprons}

%{\middel}
%&\mbox{()}\\\hline
\begin{somme}
昨年&past + year = \textit{last year}&サク{\middel}ネン \mbox{(SAKU{\middel}NEN)}\\\hline
昨夜&past + night = \textit{last night}&サク{\middel}ヤ \mbox{(SAKU{\middel}YA)}\\\hline
昨夕&past + evening = \textit{last evening}&サク{\middel}ゆう \mbox{(SAKU{\middel}Y\={U})}\\\hline
昨年来&last year + come = \textit{since last year}&サク{\middel}ネン{\middel}ライ \mbox{(SAKU{\middel}NEN{\middel}RAI)}\\\hline
昨年度&last year + degree = \textit{previous year (fiscal, academic)}&サク{\middel}ネン{\middel}ド \mbox{(SAKU{\middel}NEN{\middel}DO)}\\\hline
\notyet{昨今}&past + now = \textit{recently; nowadays}&サッ{\middel}コン \mbox{(SAK{\middel}KON)}\\\hline
\end{somme}

\begin{irrr}
昨日&past + sun (day) = \textit{yesterday}&きのう \mbox{(kin\={o})}\\\hline
\rye{3}{This compound can also be read as サク{\middel}ジツ, giving it a more formal tone.}
\end{irrr}

\begin{decomp}{3}
昨日&は&寒かった。\\\hline
きのう&は&さむかった\\\hline
yesterday&は&cold\\\hline
\multicolumn{3}{|r|}{\textit{It was cold yesterday.}}\\\hline
\end{decomp}

\end{multicols}
%%--------------------------------------------------------------------
%26-------------------------------------------------------------------
\begin{multicols}{2}
\begin{glyph}{Bundle (束)}{bundle}\label{glyph:bundle}
Bundle.  Bunch.\\\hline
The stroke order of this 漢字 is similar to 東 from Entry \ref{glyph:east} (page \pageref{glyph:east}).
\end{glyph}
\gindex{Bundle}{束}\strindex{7}{束}

\begin{comps}
\comprtwocone&木 + 口\\\hline
\comprthreecone&Tree + Mouth\\\hline
\end{comps}

\begin{remglyph}
\hooglig{Bunch} of \remcom{vampires} forming \hooglig{bundles} of \remcom{wood} from the \remcom{tree}.\\\hline
\end{remglyph}

\begin{prons}
\pronsrtwocone&\textbf{ソク (SOKU)}&\textbf{たば (taba)}\\\hline
\pronsrthreecone&---&つか (tsuka)\\\hline
\end{prons}

\begin{remprons}
\hooglig{A bundle} of \ucomkun{tsukas} were sent to \comkun{Thaba'NChu} in a \comon{sock}.\\\hline
{\warn}\textit{Thaba'NChu} is a town in Free State, South Africa.\\\hline
\end{remprons}

%{\middel}
%&\mbox{()}\\\hline
\begin{somme}
束&\textit{bundle; bunch; sheaf}&たば \mbox{(taba)}\\\hline
花束&flower + bundle = \textit{bouquet}&はな{\middel}たば \mbox{(hana{\middel}taba)}\\\hline
結束&bind + bundle = \textit{union; combination}&ケッ{\middel}ソク \mbox{(KES{\middel}SOKU)}\\\hline
束の間&bundle + interval = \textit{brief moment}&つか{\middel}の{\middel}ま \mbox{(tsuka{\middel}no{\middel}ma)}\\\hline
結束&bind + bundle = \textit{union; solidarity}&ケッ{\middel}ソク \mbox{(KES{\middel}SOKU)}\\\hline
不束&not + bundle = \textit{rude; inexperienced; stupid}&ふつ{\middel}つか \mbox{(futsu{\middel}tsuka)}\\\hline
集束&gather + bundle = \textit{focusing; convergence}&シュウ{\middel}ソク \mbox{(SH\={U}{\middel}SOKU)}\\\hline
毛束&fur + bundle = \textit{tuft of hair}&け{\middel}たば \mbox{(ke{\middel}taba)}\\\hline
\end{somme}

\begin{decomp}{3}
約束&は&約束。\\\hline
ヤクソク&は&ヤクソク\\\hline
promise&は&promise\\\hline
\multicolumn{3}{|r|}{\textit{A promise is a promise.}}\\\hline
\end{decomp}

\end{multicols}
\pagebreak
%%--------------------------------------------------------------------
%27-------------------------------------------------------------------
\begin{multicols}{2}
\begin{glyph}{Special (特)}{special}\label{glyph:special}
Special.
\end{glyph}
\gindex{Special}{特}\strindex{10}{特}

\begin{comps}
\comprtwocone&牛 + 寺\\\hline
\comprthreecone&Cow + Temple\\\hline
\end{comps}

\begin{remglyph}
\hooglig{Special} \remcom{cows} for a \remcom{temple} service.\\\hline
\end{remglyph}

\begin{prons}
\pronsrtwocone&\textbf{トク (TOKU)}&---\\\hline
\pronsrthreecone&---&---\\\hline
\end{prons}

\begin{remprons}
\hooglig{Special} \comon{talk}-shows.\\\hline
\end{remprons}

%{\middel}
%&\mbox{()}\\\hline
\begin{somme}
特別&special + different = \textit{special; extraordinary}&トク{\middel}ベツ \mbox{(TOKU{\middel}BETSU)}\\\hline
特売&special + sell = \textit{bargain; special buy}&トク{\middel}バイ \mbox{(TOKU{\middel}BAI)}\\\hline
特急&special + hasten = \textit{special express train}&トッ{\middel}キュウ \mbox{(TOKU{\middel}KY\={U})}\\\hline
特定&special + fixed = \textit{specific; particular}&トク{\middel}テイ \mbox{(TOKU{\middel}TEI)}\\\hline
特色&special + colour = \textit{characteristic; feature}&トク{\middel}ショク \mbox{(TOKU{\middel}SHOKU)}\\\hline
{\diff}特大&special + large = \textit{extra-large; king-size}&トク{\middel}ダイ \mbox{(TOKU{\middel}DAI)}\\\hline
特集&special + gather = \textit{special edition}&トク{\middel}シュウ \mbox{(TOKU{\middel}SH\={U})}\\\hline
\end{somme}

\begin{decomp}{4}
特別&いい&気分&だ。\\\hline
トクベツ&いい&キブン&だ\\\hline
special&good&feeling&be / is\\\hline
\multicolumn{4}{|r|}{\textit{I feel just fine.}}\\\hline
\end{decomp}

\end{multicols}
%%--------------------------------------------------------------------
%28-------------------------------------------------------------------
\begin{multicols}{2}
\begin{glyph}{True (真)}{true}\label{glyph:true}
True.  Genuine.\\\hline
This 漢字 often intensifies the quality of the character that follows it.  The irregular readings and the first compound below offer examples of this.
\end{glyph}
\gindex{True}{真}\strindex{10}{真}

\begin{comps}
\comprtwocone&十 + 具\\\hline
\comprthreecone&Ten + Tool (\ref{glyph:tool})\\\hline
\end{comps}

\begin{remglyph}
\hooglig{True} \remcom{cross} \remcom{tool}.\\\hline
\end{remglyph}

\begin{prons}
\pronsrtwocone&\textbf{シン (SHIN)}&\textbf{ま (ma)}\\\hline
\pronsrthreecone&---&---\\\hline
\end{prons}

\begin{remprons}
\hooglig{Truth} \comon{sheens} through the \comkun{mucky}, \comkun{muddy} lies.\\\hline
\end{remprons}

%{\middel}
%&\mbox{()}\\\hline
\begin{somme}
真新しい&true + new = \textit{brand-new}&ま　あたら{\middel}しい \mbox{(ma atara{\middel}sh\={i})}\\\hline
真実&true + actual = \textit{truth; reality}&シン{\middel}ジツ \mbox{(SHIN{\middel}JITSU)}\\\hline
真空&true + empty = \textit{vacuum; empty}&シン{\middel}クウ \mbox{(SHIN{\middel}K\={U})}\\\hline
真面目&true + face = \textit{serious; honest}&ま{\middel}じ{\middel}め \mbox{(ma{\middel}ji{\middel}me)}\\\hline
\end{somme}

\begin{irrr}
真赤&true + red = \textit{deep red}&まっ{\middel}か \mbox{(mak{\middel}ka)}\\\hline
真青&true + blue = \textit{deep blue}&まっ{\middel}さお \mbox{(mas{\middel}sao)}\\\hline
真白&true + white = \textit{snow white}&まっ{\middel}しろ \mbox{(mas{\middel}shiro)}\\\hline
真黒&true + black = \textit{pitch black}&まっ{\middel}くろ \mbox{(mak{\middel}kuro)}\\\hline
真暗&true + dark = \textit{pitch dark}&まっ{\middel}くら \mbox{(mak{\middel}kura)}\\\hline
真中&true + middle = \textit{right in the middle}&まん{\middel}なか \mbox{(man{\middel}naka)}\\\hline
\end{irrr}

\begin{decomp}{5}
この&物語&は&真実&です。\\\hline
この&ものがたり&は&シンジツ&です\\\hline
this&story&は&truth&be/is\\\hline
\multicolumn{5}{|r|}{\textit{This story is true.}}\\\hline
\end{decomp}

\end{multicols}
\pagebreak
%%--------------------------------------------------------------------
%29-------------------------------------------------------------------

\begin{multicols}{2}
\begin{glyph}{Office (局)}{office}\label{glyph:office}
Office.  Bureau.\\\hline
In a few words this 漢字 can also mean \textit{the situation}, or \textit{state of something} (the first compound is an example).
\end{glyph}
\gindex{Office}{局}\strindex{7}{局}

\begin{comps}
\comprtwocone&尸 + 口 + 勹\\\hline
\comprthreecone&Corpse/Body + Mouth + Embrace/Wrap\\\hline
\end{comps}

\begin{remglyph}
\hooglig{The bureau} of \remcom{corpses} \remcom{wrapped} the \remcom{vampires}.\\\hline
\end{remglyph}

\begin{prons}
\pronsrtwocone&\textbf{キョク (KYOKU)}&---\\\hline
\pronsrthreecone&---&---\\\hline
\end{prons}

\begin{remprons}
\hooglig{Office} of \comon{Kyokushin} karate.\\\hline
{\warn}\textit{Kyokushin} karate was founded in 1964 by Korean-Japanese Masutatsu Oyama.\\\hline
\end{remprons}

%{\middel}
%&\mbox{()}\\\hline
\begin{somme}
結局&bind + office (situation) = \textit{after all; in the long run}&ケッ{\middel}キョク \mbox{(KEK{\middel}KYOKU)}\\\hline
水道局&water supply system + office = \textit{water bureau}&スイ{\middel}ドウ{\middel}キョク \mbox{(SUI{\middel}D\={O}{\middel}KYOKU)}\\\hline
局地的&locality + target = \textit{localized; regional; isolated}&キョクチテキ \mbox{(KYOKU{\middel}CHI{\middel}TEKI)}\\\hline
局名&office + name = \textit{station name}&キョク{\middel}メイ \mbox{(KYOKU{\middel}MEI)}\\\hline
局面一転&state of the sitution + reversal = \textit{situation taking a new turn}&キョク{\middel}メン{\middel}イッ{\middel}テン (KYOKU{\middel}MEN{\middel}IT{\middel}TEN)\\\hline
局外者&the outside + individual = \textit{outsider}&キョク{\middel}ガイ{\middel}シャ \mbox{(KYOKU{\middel}GAI{\middel}SHA)}\\\hline
\notyet{局員}&office + member = \textit{staff}&キョク{\middel}イン \mbox{(KYOKU{\middel}IN)}\\\hline
\notyet{薬局}&medicine + office = \textit{pharmacy}&ヤッ{\middel}キョク \mbox{(YAK{\middel}KYOKU)}\\\hline
\notyet{局長}&office + long = \textit{office chief}&キョク{\middel}チョウ \mbox{(KYOKU{\middel}CH\={O})}\\\hline
\end{somme}
\end{multicols}

\begin{decomp}{7}
一番近い&薬局&は&どこ&に&あります&か。\\\hline
イチバンちかい&ヤッキョク&は&どこ&に&あります&か\\\hline
closest&pharmacy&は&where&に&to be&か\\\hline
\multicolumn{7}{|r|}{\textit{Where's the nearest drugstore?}}\\\hline
\end{decomp}

\pagebreak
%%--------------------------------------------------------------------
%30-------------------------------------------------------------------

\begin{multicols}{2}
\begin{glyph}{Express (表)}{express}\label{glyph:express}
Express one's feelings and beliefs, etc.\\\hline
This idea of bringing things into the open extends to this 漢字's other primary meanings, namely \textit{surfaces} and \textit{charts}.\\\hline
In some respects, this character is similar to the 漢字 for \textit{Present}, 現 (Entry \ref{glyph:present}, page \pageref{glyph:present}).  However, 表 is used more for the idea of expressing things such as feelings, thoughts, and opinions.
\end{glyph}
\gindex{Express}{表}\strindex{8}{表}

\begin{comps}
\comprtwocone&主 + 衣衤\\\hline
\comprthreecone&Primary \ref{glyph:primary} + Clothing\\\hline
\end{comps}

\begin{remglyph}
\hooglig{Express} your feelings about the \remcom{primary} school's \remcom{clothing}.\\\hline
\end{remglyph}

\begin{prons}
\pronsrtwocone&\textbf{ヒョウ (HY\={O})}&\textbf{おもて、あらわ (omote, arawa)}\\\hline
\rye{3}{ヒョウ is the more common of these readings.\\\hline
あらわ forms the two verbs in the examples below.  They are pronounced identically to 現れる and 現す.  Note that this pair can also interchangeably be written as 表れる and 表す.}
\pronsrthreecone&---&---\\\hline
\end{prons}

\begin{remprons}
\hooglig{Expressions} of \comon{yawning} when \comkun{commotions} go on in the \comkun{Arawaks}.\\\hline
{\warn}\textit{Arawaks} are indigenous people of South America.\\\hline
\end{remprons}

%{\middel}
%&\mbox{()}\\\hline
\begin{somme}
表&\textit{surface; front}&おもて \mbox{(omote)}\\\hline
表れる (intr.)&\textit{be expressed}&あらわ{\middel}れる \mbox{(arawa{\middel}reru)}\\\hline
表す (tr.)&\textit{express; reveal}&あらわ{\middel}す \mbox{(arawa{\middel}su)}\\\hline
表現&express + present = \textit{expression; representation}&ヒョウ{\middel}ゲン \mbox{(HY\={O}{\middel}GEN)}\\\hline
表現の自由&expression + freedom = \textit{freedom of expression}&ヒョウ{\middel}ゲンの{\middel}ジ{\middel}ユウ (HY\={O}{\middel}GEN{\middel}no{\middel}JI{\middel}Y\={U})\\\hline
表紙&express + paper = \textit{book cover; binding}&ヒョウ{\middel}シ \mbox{(HY\={O}{\middel}SHI)}\\\hline
表記&express + note = \textit{expression in writing; written representation}&ヒョウ{\middel}キ \mbox{(HY\={O}{\middel}KI)}\\\hline
表面的&express + mask + target = \textit{superficial; cosmetic}&ヒョウ{\middel}メン{\middel}テキ \mbox{(HY\={O}{\middel}MEN{\middel}TEKI)}\\\hline
表口&express + mouth = \textit{front door}&おもて{\middel}ぐち \mbox{(omote{\middel}guchi)}\\\hline
表皮&surface + skin = \textit{epidermis; cuticle}&ホウ{\middel}ヒ \mbox{(H\={O}{\middel}HI)}\\\hline
表意文字&ideography + character = \textit{ideogram; ideograph}&ヒョウ{\middel}イ{\middel}モ{\middel}ジ \mbox{(HY\={O}{\middel}I{\middel}MO{\middel}JI)}\\\hline
表通り&surface + passage = \textit{main street}&おもて{\middel}どお{\middel}り \mbox{(omote{\middel}d\={o}{\middel}ri)}\\\hline
表題&express + topic = \textit{title; heading}&ヒョウ{\middel}ダイ \mbox{(HY\={O}{\middel}DAI)}\\\hline
\end{somme}
\end{multicols}

\begin{decomp}{5}
思想&は&言葉&で&表現される。\\\hline
シソウ&は&ことば&で&ヒョウゲンさせる\\\hline
thought&は&words&で&expression\\\hline
\multicolumn{5}{|r|}{\textit{Thoughts are expressed by means of words.}}\\\hline
\end{decomp}

\pagebreak
%%--------------------------------------------------------------------
%---------------------------------------------------------------------